%% 00_abstract.tex — Título, autoría, resumen y palabras clave.

\title{Eficiencia del análisis dinámico de seguridad asíncrono:\\
planificación de \emph{payloads} y confirmación en ejecución\\
para la reducción de falsos positivos}

\author{%
\IEEEauthorblockN{Luis Ernesto Galdámez Vides}
\IEEEauthorblockA{%
Universidad Evangélica de El Salvador (UEES)\\
San Salvador, El Salvador\\
\texttt{2026010235@cvirtualuees.edu.sv}%
}
}

\maketitle

\begin{abstract}
Los escáneres de seguridad de aplicaciones web basados en pruebas dinámicas
(\emph{Dynamic Application Security Testing}, DAST) son ampliamente utilizados
en las canalizaciones de integración continua, pero su adopción práctica se ve
limitada por tres factores: la latencia acumulada de miles de peticiones
secuenciales, el ruido generado por un elevado número de falsos positivos y la
explosión combinatoria del espacio de \emph{payloads} por punto de inyección.
Este trabajo presenta y evalúa \herramienta, un motor DAST construido
íntegramente sobre \code{asyncio} que introduce (i)~un \emph{pipeline} de
planificación de \emph{payloads} --- entrelazado por contexto, priorización por
confianza \emph{a priori} y ordenación configurable --- y (ii)~una etapa de
confirmación en tiempo de ejecución que reetiqueta cada hallazgo con un nivel
de confianza calibrado antes de reportarlo. La hipótesis central es que ambas
técnicas mejoran la relación entre detección y coste (número de peticiones)
sin incrementar la tasa de falsos positivos. Para contrastarla se diseñó un
banco de pruebas reproducible sobre OWASP Benchmark restringido a
\num{338}~objetivos de parámetro \code{GET} (\num{194}~instancias vulnerables
y \num{144}~trampas conocidas), un estudio de ablación bajo un diseño de
bloques completos al azar (RCBD) con cuatro presupuestos de peticiones y
cinco configuraciones, y una batería de pruebas no paramétricas (Friedman,
post-hoc de Nemenyi, Wilcoxon pareado con corrección de Benjamini--Hochberg y
$\delta$ de Cliff como tamaño del efecto). La reproducibilidad se garantiza
mediante un manifiesto por ejecución que fija el \emph{hash} SHA-256 del corpus
de \emph{payloads}, el \emph{commit} de \code{git} y la semilla global del
generador de números aleatorios. El \emph{pipeline} completo alcanza una
exhaustividad del \num{64}\,\% con una tasa de falsos positivos del
\num{1.4}\,\%, y satura ya con \num{20}~peticiones por parámetro. La barrida de
ablación resulta en su mayor parte nula (Friedman $p \ll 0{,}001$ pero
$\lvert\delta\rvert < 0{,}02$). Los hallazgos con valor práctico son de
\emph{economía de recursos}: el $F_1$ es máximo en un presupuesto ajustado de
\num{20}~peticiones por parámetro; ampliarlo no detecta ninguna vulnerabilidad
nueva pero degrada la precisión unos \num{5}~puntos al acumular falsos
positivos; y la priorización por confianza —única etapa con efecto medible—
solo preserva la selectividad bajo ese régimen de escasez. El techo de
detección (\num{125} de \num{194} instancias) es una limitación de cobertura
del corpus de \emph{payloads}, no del planificador.
\end{abstract}

\begin{IEEEkeywords}
seguridad de aplicaciones web, DAST, análisis dinámico, programación
asíncrona, falsos positivos, planificación de \emph{payloads}, evaluación
reproducible, OWASP Benchmark.
\end{IEEEkeywords}
